\section[Complex Determinants]{Real Parts of Complex Determinants in All Orders}
\label{sec:complex-determinants}

The determinant has the gradient table of \cref{sec:spectral-positivity},
but \(\Delta E_{m-1}=16E_{m-2}\) excludes \(\sqrt S\) in orders
at least three; see \eqref{eq:spectral-balanced-defect}.
We prove \cref{main:complex-determinants} using the product of proper
minor norms and two coefficient inequalities uniform in matrix order.
We use the Frobenius Euclidean metric
\(\langle U,V\rangle=\operatorname{Re}\tr(UV^*)\) on
\(M_m(\mathbb C)\), so that \(|U|^2=\sum_{i,j}|U_{ij}|^2\), and write
\[
 f(Z)=\det_{\mathbb C}Z,\qquad F=\operatorname{Re}f,
 \qquad C_m=\{F=0\},\qquad E^-_m=\{F<0\}.
\]
For $0\leq k\leq m$, let
\begin{equation}\label{det:minor-norms}
 E_0=1,\qquad
 E_k(Z)=\sum_{|I|=|J|=k}|\det Z_{I,J}|^2.
\end{equation}
Let $x_1,\ldots,x_m\ge0$ be the squared singular values of $Z$,
and retain the elementary symmetric polynomials $e_k$ from
\cref{sec:spectral-positivity}. Cauchy--Binet gives
\[
 E_k(Z)=\sum_{|I|=k}\det(ZZ^*)_{I,I},\qquad
 \sum_{k=0}^m E_k(Z)t^k
 =\det(\Id+tZZ^*)=\prod_{i=1}^m(1+tx_i).
\]
Comparing coefficients yields $E_k=e_k(x_1,\ldots,x_m)$;
in particular, $E_1=|Z|^2$ and $E_m=|f|^2$.
The minor invariants $E_1,\ldots,E_{m-1}$ are positive at rank
at least $m-1$. Their product vanishes precisely at rank at most
$m-2$, the singular set of the cone, and is positive along its
regular part.
We take the product candidate
\begin{equation}\label{det:potential}
 a=\left(\prod_{k=1}^{m-1}E_k\right)^{1/2},\qquad G=aF.
\end{equation}
Logarithmic derivatives add over the minor orders.
\Cref{det:divergence-formula} reduces the defect sign to a boundary
coefficient controlled by Newton inequalities and a second coefficient
controlled by a covariance decomposition.
The homogeneous degree of $G$ is $m(m+1)/2$.
The case $m=1$ is a hyperplane; the coefficient calculations assume $m\geq2$.

For $m=2,3$, the potential is
\[
 G_2=F_2E_1^{1/2},\qquad G_3=F_3(E_1E_2)^{1/2}.
\]
The first multiplier is the radius $E_1^{1/2}=|Z|$; the second also
includes the norm $E_2^{1/2}$ of the second-order minors and agrees
with the \(d=2\) Jordan potential in \cref{thm:exc-jordan-minimizing}.
The low-order geometry and determinant phase are discussed in
\cref{det:low-orders-phase}; the coefficient signs are proved in
\cref{det:first-section,det:second-section}.

\subsection{Invariant calculus and the multiplier identity}

\begin{lemma}\label{det:basic-identities}
For $1\leq j,k\leq m$, one has
\begin{align}
 \Delta F&=0,&\Gamma(F,F)&=E_{m-1},\label{det:basic-one}\\
 \Gamma(F,E_k)&=2(m-k+1)E_{k-1}F,
 &\Delta E_k&=4(m-k+1)^2E_{k-1},\label{det:basic-two}\\
 \Gamma(E_j,E_k)
 &=4\sum_{i=1}^m x_i
       \frac{\partial e_j}{\partial x_i}
       \frac{\partial e_k}{\partial x_i}.
 \label{det:basic-three}
\end{align}
\end{lemma}

\begin{proof}
The real part of a holomorphic polynomial is harmonic.  The complex
first derivatives of $f$ are its cofactors, so their squared norms sum
to $E_{m-1}$.  This proves \eqref{det:basic-one}.

For the mixed identity, use complex matrix coordinates $z_\alpha$ and
observe that
\[
 \Gamma(F,E_k)
 =2\operatorname{Re}\sum_\alpha
       \frac{\partial f}{\partial z_\alpha}
       \frac{\partial E_k}{\partial\overline z_\alpha}.
\]
The complex expression in this sum transforms by the same determinant
phase as $f$ under the left and right unitary actions.  Thus singular
value decomposition reduces the identity to a diagonal matrix, with its
determinant phase retained.  At such a matrix,
\[
 \sum_\alpha
       \frac{\partial f}{\partial z_\alpha}
       \frac{\partial E_k}{\partial\overline z_\alpha}
 =f\sum_i e_{k-1}(x_1,\ldots,\widehat x_i,\ldots,x_m)
 =(m-k+1)E_{k-1}f.
\]
The polynomial identity extends to singular matrices by continuity.

For a holomorphic polynomial $q$,
$\Delta|q|^2=4\sum_\alpha|\partial q/\partial z_\alpha|^2$.
Apply this to the minors in \eqref{det:minor-norms}.  Each fixed
$(k-1)$-minor occurs $(m-k+1)^2$ times, which gives the Laplacian
formula.

The same unitary invariance reduces \eqref{det:basic-three} to
$D=\operatorname{diag}(\sigma_1,\ldots,\sigma_m)$, where
$x_i=\sigma_i^2$. At $D$, nonprincipal minors vanish, and
\[
 \mathrm d(\det Z_{I,I})_D[H]
 =\sum_{i\in I}
       \left(\prod_{\ell\in I\setminus\{i\}}\sigma_\ell\right)H_{ii}.
\]
Differentiating their squared moduli and summing gives
\[
 \mathrm dE_k(D)[H]
 =2\sum_i\sigma_i\frac{\partial e_k}{\partial x_i}
                                     \operatorname{Re}H_{ii}.
\]
Thus all other first derivatives vanish and
\[
 \nabla E_k(D)=\operatorname{diag}\left(
             2\sigma_i\frac{\partial e_k}{\partial x_i}\right)_{i=1}^m.
\]
Pairing these gradients proves \eqref{det:basic-three}, including
repeated or zero singular values.
\end{proof}

Apply \cref{prop:spectral-smallest-algebra} with \(d=m\), roots
\(x_i\), and \(b_k=4(m-k+1)^2>0\) from
\cref{det:basic-identities}. The algebra is Laplacian and, for \(m\ge3\),
\begin{equation}\label{det:smallest-algebra}
 \cA_R(F)=\R[E_1,\ldots,E_m,F].
\end{equation}
For \(m=2\), the same proposition gives
\(\cA_R(F)=\R[R,F]\).

The joint image of \((E_1,\ldots,E_m,F)\) is exactly
\begin{equation}\label{det:quotient-image}
 \left\{(e_1(x),\ldots,e_m(x),t):
       x\in[0,\infty)^m,\quad t\in\R,\quad t^2\le e_m(x)\right\}.
\end{equation}
Indeed, \(F^2\le|\det Z|^2=E_m\), and the diagonal matrices
\[
 Z=\operatorname{diag}\left(\sqrt{x_1}e^{i\vartheta},
                   \sqrt{x_2},\ldots,\sqrt{x_m}\right)
\]
have \(E_k=e_k(x)\) and \(F=\sqrt{e_m(x)}\cos\vartheta\).
The phase and Vandermonde argument of
\cref{prop:norm-product-exact-algebra} therefore gives the full
image, including zero spectra, and generic rank \(m+1\).
Thus for \(m\ge3\) the generators are independent and the spherical
quotient dimension is \(m\); for \(m=2\), the algebra \(\R[R,F]\)
has spherical quotient dimension one.

The identities \eqref{det:basic-one}--\eqref{det:basic-two} also give
\begin{equation}\label{det:mean-curvature}
 \cL F
 =-\tfrac12\Gamma(F,E_{m-1})=-2E_{m-2}F.
\end{equation}
Thus the regular part of $C_m$ is minimal, as in
\cite[Corollary~5.2 and Example~5.4]{HoppeTkachev19}.

\begin{example}[Generator degrees for harmonic determinant cones]
\label{ex:determinant-generator-degrees}
For \(m\ge3\), the polynomial \(F=\operatorname{Re}\det Z\) is
an irreducible homogeneous harmonic equation of degree \(m\) on
\(\R^{2m^2}\), with nonempty regular minimal zero set. Its
Laplacian algebra has \(m+1\) algebraically independent generators,
and the least possible maximal degree of homogeneous generators
is \(2m\).

Harmonicity and minimality follow from \eqref{det:basic-one} and
\eqref{det:mean-curvature}. For irreducibility, complexify the real
coordinates to independent matrices \(Z,W\), so that
\(2F=\det Z+\det W\). As a polynomial in \(z_{11}\), this is
\(Dz_{11}+C\), where \(D\) is the corresponding cofactor. Any
factor of \(D\) is independent of \(W\), whereas \(C\) has
coefficient \(1\) on \(w_{11}\cdots w_{mm}\). Thus \(D\) and
\(C\) are coprime, so \(Dz_{11}+C\) is primitive and linear over
the fraction field. Gauss' lemma proves irreducibility
\cite[Chapter~VIII, Section~5, pp.~395--398]{Knapp16}.
The zero \(Z=\operatorname{diag}(i,1,\ldots,1)\) is regular by
\eqref{det:basic-one}.

By \eqref{det:smallest-algebra},
\[
 \cA_R(F)=\R[E_1,\ldots,E_m,F],\qquad \deg E_k=2k.
\]
Algebraic independence gives the lower bound \(m+1\) for the
number of generators. Every homogeneous element of
degree less than \(2m\) belongs to the proper subalgebra
\(\R[E_1,\ldots,E_{m-1},F]\), and such elements cannot generate
\(E_m\). Conversely, the displayed generators have maximal
degree \(2m\). Thus the least possible maximal degree is exactly
\(2m\), and the uniform bound in
\cref{prop:uniform-finite-generation} must satisfy
\(B(2m^2,m)\ge2m\).
\end{example}

The singular set has
real codimension eight and is given by
\begin{equation}\label{det:singular-set}
 \mathcal S_m=\{Z:\operatorname{rank}_{\mathbb C}Z\leq m-2\}
 =\{\nabla F=0\}.
\end{equation}
The cofactor identity gives the second equality. To see that every
point of this set is geometrically singular, take a matrix of rank
$q\leq m-2$. Schur-complement coordinates, absorbing a nonvanishing
determinant factor into one row, identify its zero set locally with
a smooth factor times $\{\operatorname{Re}\det W=0\}$ at $W=0$,
where $W\in M_{m-q}(\mathbb C)$. Since $m-q\geq2$, this transverse
cone contains every complex rank-one line. Their real tangent
directions span the whole matrix space, so the zero set cannot be
a smooth hypersurface there.
The rank-\(q\) stratum has real dimension
\(2q(2m-q)\), so the largest stratum in \(\mathcal S_m\)
has dimension \(2m^2-8\).
The multiplier in \eqref{det:potential} is smooth and strictly positive
on $M_m(\mathbb C)\setminus\mathcal S_m$: at every matrix of rank at
least $m-1$, all $E_1,\ldots,E_{m-1}$ are positive. This includes the
regular zeros of $F$, even when the matrix is not invertible.
On $\mathcal S_m$ the displayed product gives $a=0$ and $G=0$.
For \(m\ge3\), this set contains nonzero points, where
\cref{prop:harmonic-positive-multiplier-obstruction} excludes multipliers
extending as strictly positive \(C^2\) functions. For \(m=2\),
\(\mathcal S_2=\{0\}\).

\begin{proposition}[Tangent cones along the determinant rank strata]
\label{prop:determinant-tangent-strata}
Let $p\in C_m$ have complex rank $q\leq m-2$, and put $r=m-q$.
The rescalings of $C_m$ about $p$ have the unique tangent cone
\begin{equation}\label{eq:determinant-rank-tangent}
 \R^{2(m^2-r^2)}\times C_r,
 \qquad C_r=\{\operatorname{Re}\det W=0:W\in M_r(\mathbb C)\},
\end{equation}
up to an orthogonal transformation. The corresponding oriented
boundary currents and their associated varifolds also converge,
with multiplicity one. The first nonzero Taylor term of $F_m$ at
$p$ has degree $r$.
\end{proposition}

\begin{proof}
The assertion at $p=0$ is homogeneity. Otherwise, unitary left and
right changes of coordinates put $p$ in the form
$\operatorname{diag}(A_0,0)$, where $A_0$ is a positive diagonal
$q$-by-$q$ matrix. These changes are Euclidean isometries; they
replace $F_m$ by $\operatorname{Re}(\chi\det Z)$ for some
$|\chi|=1$. Write nearby matrices as
\[
 Z=\begin{pmatrix}A&B\\ C&D\end{pmatrix},\qquad
 V=D-CA^{-1}B.
\]
Block elimination, with determinant one, gives
$\det Z=(\det A)\det V$. Put
\[
 \zeta(A)=\frac{\chi\det A}{|\det A|},\qquad
 U(A)=\operatorname{diag}(\zeta(A),1,\ldots,1),\qquad
 W=U(A)V.
\]
Then
\[
 \operatorname{Re}(\chi\det Z)=|\det A|\operatorname{Re}\det W.
\]
Thus the smooth local coordinate map \(\Psi(Z)=(A-A_0,B,C,W)\)
identifies both the zero set and its negative phase with a product
of a smooth coordinate space and the corresponding determinant
zero set or phase. At $p$, the differential of the Schur complement
is $dV=dD$, and the derivative of $U(A)$ contributes nothing
because $V(p)=0$. Consequently $D\Psi_p$ is the identity on the
free blocks and the unitary row transformation $D\mapsto U(A_0)D$
on the last block. In particular, it is an isometry.

For $s\downarrow0$, the maps $s^{-1}\Psi(p+sH)$ converge in
$C^1$ on compact sets to $D\Psi_p$. Since the target zero set and
phase are products of cones, this proves local Hausdorff and
local $L^1$ convergence of their inverse images. The same local
diffeomorphisms push forward the multiplicity-one boundary
current. Their $C^1$ convergence and the area formula give local
current and mass convergence, hence convergence of the associated
varifolds. This also proves uniqueness. Finally, the determinant
identity gives the leading term
\[
 \operatorname{Re}\bigl(\chi\det(A_0)\det H_{22}\bigr),
\]
which has degree $r$ and the zero set in
\eqref{eq:determinant-rank-tangent}.
\end{proof}

\begin{example}[Cubic transverse tangents at nonvertex points]
\label{ex:cubic-nonvertex-tangents}
For $F_4=\operatorname{Re}\det Z$ on $M_4(\mathbb C)$, take
$p=\operatorname{diag}(1,0,0,0)$. Then
\[
 \operatorname{Tan}_p C_4=\R^{14}\times
 \{\operatorname{Re}\det W=0:W\in M_3(\mathbb C)\}.
\]
The transverse equation is cubic. In contrast, the rank-two stratum
has quadratic transverse tangents. Both the original cone and the
cubic factor are area minimizing by \cref{main:complex-determinants}.
Their regular minimality is the determinant construction of
\cite[Corollary~5.2 and Example~5.4]{HoppeTkachev19}.

The normed-product family gives an equally direct example. Let
$P_d(x_1,\ldots,x_d)$ denote the left-associated octonionic product,
and let $\mathcal O_d=\{\operatorname{Re}P_d=0\}$. At
$p=(1,0,0,0)\in\mathcal O_4$, the coordinate change
\[
 (a,b,c,d)\longmapsto(a,(1+a)b,c,d)
\]
is a local diffeomorphism whose derivative at zero is the identity.
It transforms $\operatorname{Re}P_4(1+a,b,c,d)$ into
$\operatorname{Re}((b'c)d)$. Hence the unique set and boundary-current
tangent is $\R^8\times\mathcal O_3$. The same argument gives
$\R^8\times\mathcal O_r$ at $(1,0,\ldots,0)\in\mathcal O_{r+1}$
for every $r\geq2$. These cones are irreducible and area minimizing
by \cref{prop:norm-product-closure,thm:octonionic-product-minimizing}.

These are not cylindrical extensions of their defining equations.
For determinants, evaluation of directional derivatives on
rank-$(m-1)$ matrices detects every matrix entry of the direction.
For normed products, the directional derivative terms have distinct
block multidegrees, and $\operatorname{Re}(uv)$ is a nondegenerate
bilinear form. In both cases $D_vF\equiv0$ forces $v=0$.
The Euclidean factor appears in the tangent at the singular point;
the radial line splitting also follows from
\cref{sec:tangent-instability-transfer}.
\end{example}


For comparison, if a homogeneous cubic $F$ is defined on its
essential space, every nonzero critical point $p$ has
$\Hess F(p)\ne0$, since $D_pF(x)=\tfrac12\Hess F(p)[x,x]$
and essentiality excludes a nonzero direction with $D_pF\equiv0$.
Thus its first nonzero Taylor term is quadratic; for the harmonic
minimal case, see \cref{prop:cubic-general-singular-spectrum,sec:tangent-instability-transfer}.
The higher-order transverse cones above occur in defining degree four and higher.


We now specialize \cref{thm:exact-multiplier} to a positive smooth
invariant multiplier $a=a(E_1,\ldots,E_m)$ on the full-rank locus.
By \eqref{det:basic-two}, the chain rule gives
\[
 \Gamma(F,a)
 =F\sum_{k=1}^m2(m-k+1)E_{k-1}
                   \frac{\partial a}{\partial E_k}.
\]
Thus division by $F$ below gives smooth invariant functions.
Put $Q=E_{m-1}$, $P=a^{-2}|\nabla G|^2$, and
\begin{equation}\label{det:divergence-notation}
 c=\frac{\Gamma(F,a)}{aF},\qquad
 t=\frac{\Gamma(a,a)}{a^2},\qquad
 K=2c+t,\qquad \delta=\frac{\Delta a}{a}+c-t.
\end{equation}
For an invariant coefficient $q=q(E_1,\ldots,E_m)$, write
\begin{equation}\label{det:T-operator}
 \mathcal Tq=\frac{\Gamma(F,q)}F
 =\sum_{k=1}^m2(m-k+1)E_{k-1}\frac{\partial q}{\partial E_k}.
\end{equation}

\begin{lemma}[The multiplier identity]\label{det:divergence-formula}
With the notation just introduced, one has
\begin{align}
 P&=Q+F^2K,\label{det:P}\\
 \cL(G)&=a^3F(D_0+F^2D_1),
 \label{det:divergence}\\
 D_0&=Q\delta-2E_{m-2}-QK-\frac{\Gamma(a,Q)}{2a},
 \label{det:D0}\\
 D_1&=K\delta-\tfrac12\left(\mathcal TK+2cK+\frac{\Gamma(a,K)}a\right).
 \label{det:D1}
\end{align}
\end{lemma}

\begin{proof}
Apply \cref{lem:harmonic-log-multiplier} with \(h=\log a\).
Its notation becomes
\[
 \begin{gathered}
 S=Q,\quad q=2E_{m-2},\quad \tau=c,\quad H_h=t,\\
 Q_h=K,\quad \Delta h=\delta-c,\quad \kappa=\mathcal TK.
 \end{gathered}
\]
The last identity follows from \eqref{det:T-operator}; all these
coefficients are smooth across \(F=0\).
Substituting in \eqref{eq:harmonic-log-gradient} and
\eqref{eq:harmonic-log-two-coefficients}, with
\(\Gamma(h,Q)=\Gamma(a,Q)/a\) and
\(\Gamma(h,K)=\Gamma(a,K)/a\), gives the stated formulas.
\end{proof}

In the notation of \cref{thm:exact-multiplier}, the lemma states
\[
 \cK_F(a)=a^3(D_0+F^2D_1),\qquad
 G\cL(G)=aF^2\cK_F(a)=a^4F^2(D_0+F^2D_1).
\]
The coefficient $D_0$ determines the boundary Jacobi sign, while
$D_1>0$ supplies strict divergence off the cone once $D_0\geq0$.
We now take the uniform product multiplier
\begin{equation}\label{det:uniform-weight}
 a=\prod_{k=1}^{m-1}E_k^\lambda.
\end{equation}
The value $\lambda=1/2$ gives \eqref{det:potential}.

\subsection{From the coefficient signs to area minimization}
\label{det:geometric-proof}

We first deduce the geometric conclusions from the coefficient signs
proved in \cref{det:first-section,det:second-section}.

\begin{proof}[Proof of \cref{main:complex-determinants}]
The case $m=1$ is an oriented hyperplane.

\emph{The sign of $G\cL(G)$.}  For $m\geq2$, take
$\lambda=1/2$ in \eqref{det:uniform-weight}.  Then
$G=aF$ is the potential in \eqref{det:potential}, and
\cref{det:first-reduction,det:first-inequality,det:second-positive}
give $D_0,D_1\geq0$.  Therefore \eqref{det:divergence} yields
\begin{equation}\label{det:divergence-signs}
 G\cL(G)=a^4F^2(D_0+F^2D_1)\geq0
 \qquad\text{on the full-rank locus}.
\end{equation}

\emph{Extension and comparison.}
Off $\mathcal S_m$, the proper minor norms are positive, so $a,G$
are smooth and \eqref{det:divergence-signs} extends by continuity to
rank $m-1$. Positive homogeneity gives $\nabla G\ne0$ there by the
verification following \cref{cor:gradient-subcalibration}.
The continuous potential has the phases of $F$, and the exceptional
set has real codimension eight by \eqref{det:singular-set}.
Thus \cref{cor:gradient-subcalibration,thm:removable-subcalibration}
give perimeter and mass comparison; harmonicity gives full support
by \cref{rem:polynomial-finite-perimeter}.

\emph{Strict minimality.}  For $\lambda=1/2$, \cref{det:second-positive} gives $D_1>0$.
Hence $q_{X_G}=\operatorname{sgn}(G)\diver X_G>0$ at full-rank points
with $F\ne0$. Each phase contains such points, giving positive flux
integrals. Homogeneity and \cref{thm:strict-minimality-flux} prove
strict area minimality.

\emph{Stability.} Off \(\mathcal S_m\), all of
\(E_1,\ldots,E_{m-1}\) are positive, so the invariant expression
\eqref{det:D0} is smooth.
Thus \(D_0\geq0\) extends from full-rank matrices to rank \(m-1\)
by continuity. On the regular cone, the Jacobi potential in
\eqref{eq:boundary-potential-algebraic} is consequently
\[
 V_G=\frac{\cK_F(a)}{a^3|\nabla F|^2}=\frac{D_0}{E_{m-1}}\geq0.
\]
Apply \cref{thm:ground-state-gap} with
$n=2m^2-1$ and $\ell=m(m+1)/2$. This gives
\eqref{eq:main-determinant-stability}:
\[
 Q_{C_m}(u)\geq
 \left(\ell-\frac n2\right)^2\int_{C_m}r^{-2}u^2
 =\frac{(m^2-m-1)^2}{4}\int_{C_m}r^{-2}u^2.\qedhere
\]
\end{proof}

\begin{remark}[Low orders and determinant phase]\label{det:low-orders-phase}
For \(m=2\), the cone \(C_2\) is orthogonally congruent to the
classical Simons cone in \(\R^4\oplus\R^4\).
Indeed, for \(Z=\left(\begin{smallmatrix}z_1&z_2\\z_3&z_4\end{smallmatrix}\right)\),
the unitary change of coordinates
\[
 w=\frac1{\sqrt2}\bigl(z_1+z_4,\ i(z_1-z_4),\ i(z_2+z_3),\ z_2-z_3\bigr)
\]
gives \(\sum_{j=1}^4w_j^2=2\det Z\).
Writing \(w=u+iv\), with \(u,v\in\R^4\), we obtain
\(2F_2=|u|^2-|v|^2\). Thus this order recovers the standard
quadratic cone.

For every real $\theta$, the isometry $Z\mapsto e^{-i\theta/m}Z$
preserves all the $E_k$ and sends $F$ to
$\operatorname{Re}(e^{-i\theta}\det_{\mathbb C}Z)$.
Consequently the phase cone
$\{\operatorname{Re}(e^{-i\theta}\det_{\mathbb C}Z)=0\}$ has the
corresponding subcalibration and is locally area minimizing for every
$m\geq1$, with the same strict minimality and stability conclusions
when $m\geq2$. In particular, the real and imaginary determinant zero cones
are orthogonally congruent.
\end{remark}

\subsection{Nonnegativity of the first coefficient}
\label{det:first-section}

We prove $D_0\geq0$ for the choice $\lambda=1/2$ in
\eqref{det:uniform-weight}.  On the full-rank locus, set
$y_i=x_i^{-1}$. For the remainder of this section, an unadorned
$e_s$ denotes the elementary symmetric polynomial in $y_1,\ldots,y_m$.
For $1\leq s\leq m-1$, define
\begin{equation}\label{det:a-variables}
 a_{si}=\frac{e_s(y_1,\ldots,\widehat y_i,\ldots,y_m)}{e_s(y)},
 \qquad P_i=\sum_{s=1}^{m-1} a_{si},
\end{equation}
and introduce four scalar sums for the coefficient calculation:
\begin{equation}\label{det:A-V-Y}
 A=\sum_i y_i\sum_{s=1}^{m-1}(s a_{si}-a_{si}^2),\qquad
 B=\sum_i y_iP_i,\qquad
 V=\sum_i y_iP_i^2,\qquad Y=\sum_i y_i^2P_i.
\end{equation}

\begin{lemma}\label{det:first-reduction}
For the weight \eqref{det:uniform-weight},
\begin{equation}\label{det:first-reduction-equation}
 \frac{D_0}{Q}
 =4\lambda A-4\lambda^2V
       -2\frac{e_2}{e_1}+2\lambda\frac{Y}{e_1}.
\end{equation}
Consequently, at $\lambda=1/2$, the inequality $D_0\geq0$ is equivalent
to
\begin{equation}\label{det:first-homogeneous}
 2A+\frac{Y}{e_1}-V-2\frac{e_2}{e_1}\geq0.
\end{equation}
\end{lemma}

\begin{proof}
We first explain the reciprocal substitution. Put $p=\prod_i x_i$.
For a $k$-element index set $I$,
$\prod_{i\in I}x_i=p\prod_{j\notin I}y_j$. Summing over $I$ gives
\[
 E_k(x)=p\,e_{m-k}(y),\qquad
 \frac{E_{k-1}(x)}{E_k(x)}
 =\frac{e_{m-k+1}(y)}{e_{m-k}(y)}.
\]
Write $y_{\widehat i}$ for the tuple with $y_i$ omitted. Since \(\partial_{y_i}e_{s+1}(y)=e_s(y_{\widehat i})\),
Euler's identity for the homogeneous polynomial $e_{s+1}$ gives
\[
 \sum_i y_i e_s(y_{\widehat i})=(s+1)e_{s+1}(y),\qquad
 \sum_i y_i a_{si}=(s+1)\frac{e_{s+1}}{e_s}.
\]
For $s=m-k$, applying the complement identity in the $m-1$ variables
with $x_i$ omitted also gives
\[
 e_{k-1}(x_{\widehat i})=p\,y_i e_s(y_{\widehat i}),\qquad
 \frac{e_{k-1}(x_{\widehat i})}{E_k(x)}=y_i a_{si}.
\]
Consequently, the gradient formula \eqref{det:basic-three} yields
\[
 \frac{\Gamma(E_k,E_k)}{E_k^2}
 =4\sum_i x_i\left(\frac{e_{k-1}(x_{\widehat i})}{E_k(x)}\right)^2
 =4\sum_i y_i a_{si}^2,
\]
where $x_i y_i^2=y_i$.
Put
$h=\log a=\lambda\sum_{k=1}^{m-1}\log E_k$.
The chain rule in \eqref{eq:log-multiplier-derivatives} gives
\[
 \frac{\nabla a}{a}=\lambda\sum_{k=1}^{m-1}\frac{\nabla E_k}{E_k},
 \qquad
 \frac{\Delta a}{a}-t
 =\lambda\sum_{k=1}^{m-1}
 \left(\frac{\Delta E_k}{E_k}
       -\frac{\Gamma(E_k,E_k)}{E_k^2}\right).
\]
Using \cref{det:basic-identities} and the identities just proved, we
obtain the quantities defined in \eqref{det:divergence-notation}:
\begin{equation}\label{det:radial-identities}
 \begin{gathered}
 c=2\lambda B,\qquad t=4\lambda^2V,\qquad
 \frac{\Delta a}{a}-t=4\lambda(A+B),\\
 \delta=4\lambda A+6\lambda B,\qquad
 K=4\lambda B+4\lambda^2V.
 \end{gathered}
\end{equation}
Finally,
\[
 \frac{E_{m-2}}Q=\frac{e_2}{e_1},\qquad
 \frac{\Gamma(a,Q)}{aQ}=4\lambda\left(B-\frac{Y}{e_1}\right).
\]
Insert these identities in \eqref{det:D0}.  The three terms involving
$B$ cancel, leaving \eqref{det:first-reduction-equation}.
\end{proof}

\begin{theorem}[First-coefficient inequality]\label{det:first-inequality}
Let $m\geq2$ and let $y_1,\ldots,y_m$ be positive.
Then the quantities in \eqref{det:a-variables}--\eqref{det:A-V-Y}
satisfy \eqref{det:first-homogeneous}.
\end{theorem}

We prove the theorem by a recurrence for symmetric polynomials.  Put
\begin{equation}\label{det:cleared-polynomial}
 D_m^{\mathrm{sym}}=\prod_{s=1}^{m-1}e_s,\qquad
 \mathcal C_m=e_1(D_m^{\mathrm{sym}})^2
       \left(2A+\frac{Y}{e_1}-V-2\frac{e_2}{e_1}\right).
\end{equation}
This is a symmetric polynomial, viewed as a polynomial in
$e_1,\ldots,e_m$. In the first two orders,
\[
 \mathcal C_2=0,\qquad
 \mathcal C_3=6e_1e_2e_3(e_1^2-2e_2).
\]
Thus the first coefficient vanishes identically for $m=2$; for
$m=3$, the normalization $e_1(y)=1$ gives
\[
 2A+Y-V-2e_2
 =6e_3\frac{\sum_i y_i^2}{e_2}>0.
\]
For $N\geq m$, let $\mathcal C_m^{[N]}$ denote
this same formal polynomial evaluated at the elementary symmetric
functions of $N$ positive variables. Its coefficients and summation
limits remain those defining \(\mathcal C_m\); they are not recomputed
by replacing \(m\) with \(N\) in \eqref{det:cleared-polynomial}.
In the induction, \(N\) is fixed and \(r\) increases, so
\(\mathcal C_r^{[N]}\) is evaluated on a fixed \(N\)-tuple.
The formal recurrence below remains valid under these evaluations.

\begin{lemma}[Recurrence for the symmetric polynomials]\label{det:stable-recurrence}
One has $\mathcal C_2=0$.  For every $r\geq2$,
\begin{equation}\label{det:C-recurrence}
 \mathcal C_{r+1}=e_r^2\mathcal C_r+\mathcal R_r.
\end{equation}
With
\[
 \mathcal P_r=e_1^3e_re_{r+1}\prod_{j=2}^{r-1}e_j^2,
\]
the remainder is given by
\begin{align}
 \frac{\mathcal R_r}{\mathcal P_r}
 ={}&r(r+1)
 -\sum_{p=1}^{\lfloor r/2\rfloor}
       (r+1)(2+\mathbf1_{p=1}+\mathbf1_{2p=r})
       \frac{e_r}{e_pe_{r-p}}\notag\\
 &-\sum_{d=1}^{\lfloor(r-1)/2\rfloor}
       \sum_{k=1}^{r-2d}
       (r-k+1)(2+\mathbf1_{k=r-2d})
       \frac{e_ke_r}{e_{k+d}e_{r-d}}.
 \label{det:R-recurrence}
\end{align}
Every denominator on the right cancels against $\mathcal P_r$.
\end{lemma}

The proof is given in \cref{app:determinant-details}.

Newton's inequalities bound the sums in \eqref{det:R-recurrence}
by the following binomial sums, whose total is at most \(r(r+1)\).
This will give \(\mathcal R_r^{[N]}\geq0\).

\Needspace{10\baselineskip}
\begin{lemma}[Scalar binomial inequality]\label{det:binomial-budget}
For $r\geq2$, set
\begin{align*}
 I_r&=\sum_{p=1}^{\lfloor r/2\rfloor}
       \frac{(r+1)(2+\mathbf1_{p=1}+\mathbf1_{2p=r})}{\binom rp},\\
 J_r&=\sum_{d=1}^{\lfloor(r-1)/2\rfloor}\sum_{k=1}^{r-2d}
       (r-k+1)(2+\mathbf1_{k=r-2d})
       \frac{\binom{k+d}{d}}{\binom rd}.
\end{align*}
Then
\begin{equation}\label{det:scalar-budget}
 I_r+J_r\leq r(r+1).
\end{equation}
\end{lemma}

The proof is given in \cref{app:determinant-details}.

\begin{proof}[Proof of \cref{det:first-inequality}]
Fix $N\geq r+1$ positive variables.  Newton's inequalities state that
the normalized sequence $\overline e_j=e_j/\binom Nj$ satisfies
\(\overline e_j^{\,2}\geq\overline e_{j-1}\overline e_{j+1}\).
This is its log concavity.  Hence
\begin{align}
 \frac{e_r}{e_pe_{r-p}}
 &\leq\frac{\binom Nr}{\binom Np\binom N{r-p}}
 \leq\frac1{\binom rp},\label{det:newton-one}\\
 \frac{e_ke_r}{e_{k+d}e_{r-d}}
 &\leq\frac{\binom Nk\binom Nr}
           {\binom N{k+d}\binom N{r-d}}
 \leq\frac{\binom{k+d}{d}}{\binom rd}.
 \label{det:newton-two}
\end{align}
For the last inequality, the quotient of the middle expression by the
last expression is
\[
 \prod_{j=1}^d\frac{N-r+j}{N-k-d+j}\leq1,
\]
since $k+d\leq r$.  Thus \cref{det:binomial-budget} and
\eqref{det:R-recurrence} give $\mathcal R_r^{[N]}\geq0$ for every
$N\geq r+1$.

We induct in $r$ with the assertion
$\mathcal C_r^{[N]}\geq0$ for every $N\geq r$.  The base is
$\mathcal C_2=0$, and \eqref{det:C-recurrence} proves the next step.
Taking $N=m$ and dividing by the positive factor in
\eqref{det:cleared-polynomial} proves
\eqref{det:first-homogeneous}.  Equivalently, under $e_1(y)=1$, this is
$2A+Y-V-2e_2\geq0$.
\end{proof}

\subsection{Nonnegativity of the second coefficient}
\label{det:second-section}

We prove $D_1\geq0$ by a covariance formula valid for every
$\lambda\geq0$, including $\lambda=1/2$. We retain the variables
$y_i=x_i^{-1}$.
\begin{lemma}\label{det:second-positive}
For the uniform weight \eqref{det:uniform-weight}, $D_1\geq0$ on the
full-rank locus whenever $\lambda\geq0$. If $\lambda>0$, then $D_1>0$.
\end{lemma}

\begin{proof}
To differentiate the ratios $a_{si}$, we use a finite probability space.
For each $1\leq s\leq m-1$, consider the probability measure on
$s$-element subsets of $\{1,\ldots,m\}$ given by
\[
 \mathbb P_s(S)=\frac{\prod_{i\in S}y_i}{e_s(y)}.
\]
This is the conditional Bernoulli distribution of
\cite[pp.~876--877, equations~(3) and~(6); Section~3.2, p.~883]{ChenLiu97}, with weights
\(y_i\) and fixed sample size \(s\).
For random variables \(\xi,\eta\) on this finite probability space,
their covariance is
\[
 \operatorname{Cov}_s(\xi,\eta)
 =\mathbb E_s(\xi\eta)-\mathbb E_s(\xi)\mathbb E_s(\eta).
\]
Let \(\xi_i(S)\) be the indicator of \(i\in S\), equal to one
when \(i\) belongs to \(S\) and zero otherwise.
Define
\[
 \pi_{si}=\mathbb P_s(i\in S)=1-a_{si},\qquad
 w_{sij}=\pi_{si}\pi_{sj}-\mathbb P_s(i,j\in S)\quad(i\ne j).
\]
Thus \(w_{sij}=-\operatorname{Cov}_s(\xi_i,\xi_j)\).
If $e_q^{ij}$ denotes the elementary symmetric polynomial with $y_i,y_j$
omitted, then
\begin{equation}\label{det:covariance-positive}
 w_{sij}=\frac{y_iy_j}{e_s^2}
       \left[(e_{s-1}^{ij})^2-e_{s-2}^{ij}e_s^{ij}\right]\geq0.
\end{equation}
This is Newton's inequality, with the convention $e_{-1}^{ij}=0$.
It is also the Rayleigh inequality for \(e_s\), expressing pairwise
negative correlation; see
\cite[Definitions~2.3 and~2.5]{BorceaBrandenLiggett09}.
Since every subset has exactly \(s\) elements,
\(\sum_i\xi_i=s\).  This gives
\begin{equation}\label{det:covariance-rowsum}
 \sum_{j\ne i}w_{sij}=\pi_{si}(1-\pi_{si})=\pi_{si}a_{si}.
\end{equation}
Indeed, the sum of the covariances of the indicator of $i\in S$ with
all the indicators is its covariance with the constant $|S|=s$.

In addition to \eqref{det:A-V-Y}, put
\[
 L=\sum_i y_i\sum_{s=1}^{m-1}s a_{si},\qquad Z_2=\sum_i y_i^2P_i^2,
 \qquad U=\sum_i y_i^2P_i^3,
\]
and define
\begin{align}
 H_s&=\sum_i y_i^2a_{si}
       -\sum_{i<j}w_{sij}(y_i-y_j)^2,
 \label{det:H-positive}\\
 \mathcal Q_1&=\sum_{s,i<j}w_{sij}
       \left[B(y_i+y_j)-(y_i-y_j)(y_iP_i-y_jP_j)\right],
 \label{det:Q1}\\
 \mathcal Q_2&=\sum_{s,i<j}w_{sij}
       \left[V(y_i+y_j)-(y_iP_i-y_jP_j)^2\right].
 \label{det:Q2}
\end{align}

We claim the decomposition
\begin{equation}\label{det:D1-positive-formula}
 D_1=\lambda R_1+\lambda^2R_2+\lambda^3R_3,
\end{equation}
where
\begin{align}
 R_1&=4\sum_sH_s,\label{det:R1}\\
 R_2&=16LB+12Z_2+16\mathcal Q_1,\label{det:R2}\\
 R_3&=16LV+8U+16\mathcal Q_2.\label{det:R3}
\end{align}
We derive it first, then prove the coefficient signs.
Write
$\mathscr D_j=y_j\partial_{y_j}$.  Differentiating the probabilities gives
\begin{equation}\label{det:probability-derivative}
 \mathscr D_j a_{si}
 =\begin{cases}
   w_{sij},&i\ne j,\\
   -\sum_{\ell\ne i}w_{si\ell},&i=j.
  \end{cases}
\end{equation}
Indeed, direct differentiation gives
\(\mathscr D_j\mathbb P_s(S)=(\xi_j(S)-\pi_{sj})\mathbb P_s(S)\).
Thus \(-\mathscr D_j a_{si}=\operatorname{Cov}_s(\xi_i,\xi_j)\),
the usual covariance identity in the logarithmic parameters
\(\log y_j\).
Consequently,
\begin{align}
 \mathscr D_jB
 &=y_jP_j+\sum_s\sum_{i\ne j}w_{sij}(y_i-y_j),
 \label{det:DB}\\
 \mathscr D_jV
 &=y_jP_j^2+2\sum_s\sum_{i\ne j}w_{sij}(y_iP_i-y_jP_j).
 \label{det:DV}
\end{align}
For any smooth invariant function $q$ written in the $y$ variables,
the change of variables gives
$\partial_{x_j}q=-y_j\mathscr D_jq$.
The product multiplier \eqref{det:uniform-weight} satisfies
$a^{-1}\partial_{x_j}a=\lambda y_jP_j$, by the complement identity
in the proof of \cref{det:first-reduction}.
Thus \eqref{det:T-operator} and the gradient pairing
\eqref{det:basic-three} give
\begin{equation}\label{det:radial-operators}
 \mathcal Tq=-2\sum_j y_j\mathscr D_jq,\qquad
 \frac{\Gamma(a,q)}a=-4\lambda\sum_j y_jP_j\mathscr D_jq.
\end{equation}
Insert \eqref{det:radial-identities} and
\eqref{det:radial-operators} in \eqref{det:D1}.  The coefficients of
$\lambda,\lambda^2,\lambda^3$ are, respectively,
\begin{align*}
 R_1&=4\sum_jy_j\mathscr D_jB,\\
 R_2&=16B(A+B)+4\sum_jy_j\mathscr D_jV
                       +8\sum_jy_jP_j\mathscr D_jB,\\
 R_3&=16V(A+B)+8\sum_jy_jP_j\mathscr D_jV.
\end{align*}
Pairing the terms with indices $i,j$ in
\eqref{det:DB}--\eqref{det:DV} gives
\begin{align*}
 \sum_jy_j\mathscr D_jB&=\sum_sH_s,\\
 \sum_jy_j\mathscr D_jV
 &=Z_2-2\sum_{s,i<j}w_{sij}(y_i-y_j)(y_iP_i-y_jP_j),\\
 \sum_jy_jP_j\mathscr D_jB
 &=Z_2-\sum_{s,i<j}w_{sij}(y_i-y_j)(y_iP_i-y_jP_j),\\
 \sum_jy_jP_j\mathscr D_jV
 &=U-2\sum_{s,i<j}w_{sij}(y_iP_i-y_jP_j)^2.
\end{align*}
The row-sum identity \eqref{det:covariance-rowsum} also gives
\begin{equation}\label{det:AplusB}
 A+B=L+\sum_{s,i<j}w_{sij}(y_i+y_j).
\end{equation}
These identities yield \eqref{det:R1}--\eqref{det:R3}.

To prove the signs, use \eqref{det:covariance-rowsum} to obtain
\[
 \sum_{i<j}w_{sij}(y_i-y_j)^2
 \leq\sum_i y_i^2\pi_{si}a_{si}
 \leq\sum_i y_i^2a_{si}.
\]
Thus $H_s\geq0$.  Also,
\[
 |(y_i-y_j)(y_iP_i-y_jP_j)|
 \leq(y_i+y_j)(y_iP_i+y_jP_j)
 \leq B(y_i+y_j),
\]
so $\mathcal Q_1\geq0$.  Cauchy--Schwarz gives
\[
 (y_iP_i-y_jP_j)^2
 \leq(y_i+y_j)(y_iP_i^2+y_jP_j^2)
 \leq V(y_i+y_j),
\]
which proves $\mathcal Q_2\geq0$.  The other terms in
\eqref{det:R2}--\eqref{det:R3} are nonnegative by definition.
Finally, $P_i>0$ and $y_i>0$ imply $Z_2>0$.
Thus $R_2\geq12Z_2$ and $D_1\geq12\lambda^2Z_2>0$ when $\lambda>0$.
\end{proof}

Despite a common gradient-pairing table, the different Laplacian
coefficients of the spectral families in
\cref{sec:norm-products-pfaffians} and the determinants
lead to the multipliers $\sqrt S$ and the product of proper minor norms,
respectively. Algebraic closure and spectral inequalities certify both
in every order. For determinants, $m+1$ independent generators when
$m\ge3$ give unbounded defining degree and generator count.
